\documentclass{article}
\usepackage{graphicx}
\usepackage[T2A]{fontenc}
\usepackage[utf8x]{inputenc}
\usepackage[english, russian]{babel}
\usepackage{csquotes}
\usepackage{hyperref}
\usepackage{amsmath}
\usepackage{biblatex}
\addbibresource{references.bib}

\title{Обзор современных нейросетевых оптимизаторов}
\author{Колпащиков Максим}
\date{March 2025}

\begin{document}

\maketitle

\section{\textbf{AdaBelief}}

\textbf{Год выхода:} 2020

\textbf{Гиперпараметры:}

\begin{itemize}
    \item \textbf{lr} - скорость обучения (значение по умолчанию: \textbf{1e-3});
    \item \textbf{betas} - коэффициенты для вычисления средних значений градиента и их квадратов (значения по умолчанию: (\textbf{0.9}, \textbf{0.999}));
    \item \textbf{epsilon} - значения в знаменателе дроби, чтобы избежать деления на ноль или близкой к нулю величине (значение по умолчанию: \textbf{1e-8});
    \item \textbf{weight\_decay} - снижение весов (L2-регуляризация) (значение по умолчанию: \textbf{0});
    \item \textbf{amsgrad} - использовать ли вариант \textbf{AMSGrad} для этого алгоритма (значение по умолчанию: \textbf{False});
    \item \textbf{weight\_decouple} - использовать ли уменьшение весов как в \textbf{AdamW} (значение по умолчанию: \textbf{False});
    \item \textbf{fixed\_decay} - используется, если для параметра \textbf{weight\_decay} установлено значение \textbf{True}. Уменьшение весов выполняется по формуле: $W_{new} = W_{old} - W_{old} \times decay$, если False: $W_{new} = W_{old} - W_{old} \times decay \times lr$ (значение по умолчанию: \textbf{False});
    \item \textbf{rectify} - выполняет обновление, аналогично \textbf{RAdam} (значение по умолчанию: \textbf{False}).
\end{itemize}

\textbf{Идея:} Интуитивный подход к \textbf{AdaBelief} заключается в том, чтобы адаптировать размер шага в соответствии с "верой" в текущее направление градиента. Рассматривая экспоненциальную скользящую среднюю (EMA) зашумленного градиента как предсказание градиента на следующем временном шаге, если наблюдаемый градиент сильно отличается от предсказания, мы не доверяем текущему наблюдению и делаем небольшой шаг; если наблюдаемый градиент близок к прогнозируемому, мы доверяем ему и делаем большой шаг вперед. Авторы утверждают достоверность данных в ходе обширных экспериментов, показывающих, что он превосходит другие методы благодаря быстрой сходимости и высокой точности при классификации изображений и языковом моделировании.

\textbf{Источник:} \cite{adabelief}

\section{\textbf{Lion}}

\textbf{Год выхода:} 2023

\textbf{Гиперпараметры:}

\begin{itemize}
    \item \textbf{lr} - скорость обучения (значение по умолчанию: \textbf{1e-4});
    \item \textbf{betas} - коэффициенты для вычисления средних значений градиента и их квадратов (значения по умолчанию: (\textbf{0.9}, \textbf{0.99}));
    \item \textbf{weight\_decay} - снижение весов (L2-регуляризация) (значение по умолчанию: \textbf{0});
    \item \textbf{use\_triton} - использовать \textbf{Triton} (значение по умолчанию: \textbf{False});
    \item \textbf{decoupled\_weight\_decay} - использовать ли независимый спад весов (значение по умолчанию: \textbf{False}).
\end{itemize}

\textbf{Идея:} Использование знака градиента и инерции для обновления весов модели. В отличие от адаптивных методов, таких как \textbf{Adam}, \textbf{Lion} применяет одинаковую величину обновления для каждого параметра, определяемую через функцию знака. Это позволяет \textbf{Lion} уменьшить потребление памяти и скорость работы, поскольку он отслеживает только моментум, а не вторые моменты градиента.

\textbf{Источник:} \cite{lion}

\section{\textbf{AdaHessian}}

\textbf{Год выхода:} 2020

\textbf{Гиперпараметры:}

\begin{itemize}
    \item \textbf{lr} - скорость обучения (значение по умолчанию: \textbf{0.15});
    \item \textbf{betas} - коэффициенты для вычисления средних значений градиента и их квадратов (значения по умолчанию: (\textbf{0.9}, \textbf{0.999}));
    \item \textbf{epsilon} - значения в знаменателе дроби, чтобы избежать деления на ноль или близкой к нулю величине (значение по умолчанию: \textbf{1e-4});
    \item \textbf{weight\_decay} - снижение весов (L2-регуляризация) (значение по умолчанию: \textbf{0});
    \item \textbf{hessian\_power} - степень гессиана (значение по умолчанию: \textbf{1});
    \item \textbf{spatial\_average\_block\_size} - Пространственный средний размер блока для одномерных тензоров (значения по умолчанию: \textbf{(-1, -1, -1, -1)}).
\end{itemize}

\textbf{Идея:} \textbf{AdaHessian} использует информацию о гессиане (вторых производных) для адаптации скорости обучения, что улучшает сходимость.

\textbf{Источник:} \cite{adahessian}

\section{\textbf{Fromage}}

\textbf{Год выхода:} 2020

\textbf{Гиперпараметры:}

\begin{itemize}
    \item \textbf{lr} - скорость обучения (значение по умолчанию: \textbf{1e-2});
    \item \textbf{p\_bound} - ограничивает оптимизацию ограниченным набором.
\end{itemize}

\textbf{Идея:} Fromage сочетает адаптивную нормализацию градиента с контролем нормы весов, что способствует стабильности и эффективности обучения.

\textbf{Источник:} \cite{fromage}

\section{\textbf{AdaBound}}

\textbf{Год выхода:} 2019

\textbf{Гиперпараметры:}

\begin{itemize}
    \item \textbf{lr} - скорость обучения \textbf{Adam} (значение по умолчанию: \textbf{1e-3});
    \item \textbf{betas} - коэффициенты для вычисления средних значений градиента и их квадратов (значения по умолчанию: (\textbf{0.9}, \textbf{0.999}));
    \item \textbf{final lr} - окончательная (\textbf{SGD}) скорость обучения (значение по умолчанию: \textbf{1e-1});
    \item \textbf{gamma} - скорость сходимости для граничной функции (значение по умолчанию: \textbf{1e-3});
    \item \textbf{epsilon} - значения в знаменателе дроби, чтобы избежать деления на ноль или близкой к нулю величине (значение по умолчанию: \textbf{1e-8});
    \item \textbf{weight\_decay} - снижение весов (L2-регуляризация) (значение по умолчанию: \textbf{0});
    \item \textbf{amsgrad} - использовать ли вариант \textbf{AMSGrad} для этого алгоритма (значение по умолчанию: \textbf{False}).
\end{itemize}

\textbf{Идея:} AdaBound адаптирует скорость обучения с динамическими границами, переходя от адаптивных методов к SGD по мере обучения, что обеспечивает быструю сходимость и хорошую обобщающую способность.

\textbf{Источники:} \cite{adabound}

\printbibliography

\end{document}
